\begin{pastexam}[label={px:ust-107-ee-ds-10}]{臺聯大電機類 107 學年度 第十題}
（問答；10 分）

(10\%) Given the following single-pair shortest-path algorithm by depth-first search,

\underline{Single-pair shortest-path}
\begin{enumerate}[label=\arabic*.]
\item Choose the smallest unvisited child of the start node
\item Choose the smallest unvisited child of that node and so on
\item Going deeper and deeper until it either reaches the goal node or a node with no children
\item The search then backtracks, returning to the most recent node with children that it has not yet visited
\item When all paths have been explored, choose the shortest path from the start to the goal
\end{enumerate}

What is the shortest path from node 0 to node 5 of the graph of Problem 八 by this algorithm? Write down each step in detail. Assume that all edges have the same weight one on the graph. Hint: the first four steps of the algorithm should be 0, 0$\to$1, 0$\to$1$\to$2, and 0$\to$1$\to$2$\to$3.

\medskip
（第八題的圖）
\begin{center}
\begin{tikzpicture}[v/.style={circle,fill=black!75,minimum size=7mm,inner sep=0pt,text=white,font=\small},>={Stealth[length=2.2mm]},thick]
  \node[v] (0) at (0,0) {0};
  \node[v] (1) at (-1.6,-1.7) {1};
  \node[v] (2) at (2.0,-1.7) {2};
  \node[v] (3) at (2.0,-3.6) {3};
  \node[v] (4) at (3.8,-3.6) {4};
  \node[v] (5) at (2.0,-5.5) {5};
  \draw[->] (0) to[out=-120,in=60] (1);
  \draw[->] (0) to[out=-55,in=125] (2);
  \draw[->] (1) -- (2);
  \draw[->] (2) -- (3);
  \draw[->] (2) to[out=-50,in=115] (4);
  \draw[->] (3) -- (4);
  \draw[->] (3) -- (5);
  \draw[->] (3) to[out=180,in=-60] (1);
  \draw[->] (1.west) .. controls (-2.7,-1.4) and (-2.8,1.4) .. (-1.2,0.85) .. controls (-0.7,0.7) and (-0.45,0.6) .. (0.120);
  \draw[->] (4.east) .. controls (5.1,-2.5) and (4.9,0.9) .. (2.6,0.85) .. controls (1.5,0.85) and (0.7,0.6) .. (0.60);
\end{tikzpicture}
\end{center}
\end{pastexam}

\begin{pastsoln}
答案：最短路徑是 $0 \to 2 \to 3 \to 5$，長度 $3$。

圖的出邊（第八題）依編號排好：$0{:}\,1, 2$；$1{:}\,0, 2$；$2{:}\,3, 4$；$3{:}\,1, 4, 5$；$4{:}\,0$；$5$ 沒有出邊。

\textbf{「unvisited」的讀法。} 第 5 步要「把所有路徑都探索過」再挑最短的，所以「未拜訪」要讀成「不在目前這條路徑上」
（避免繞圈），回溯時要把頂點從路徑上拿掉，之後別的路徑還可以再經過它。 這其實是回溯法列出所有簡單路徑
（資結書第~20.1~節）。 若照一般 DFS 用全域的拜訪標記，走完 $0 \to 1 \to 2 \to 3 \to 5$ 之後，$2$ 已經標記過，
$0 \to 2$ 這條路就不會再走，只會找到長度 $4$ 的路徑，不是最短的；所以題目的演算法不能這樣讀。 另外，碰到目標
$5$ 就停在這條路徑、回溯（第 3 步），不再從 $5$ 往下走。

\textbf{逐步追蹤}（每一步寫出目前的路徑）：
\begin{center}
\small
\begin{tabular}{@{}rlp{0.48\linewidth}@{}}
\toprule
步 & 目前的路徑 & 說明 \\
\midrule
1 & $0$ & 起點 \\
2 & $0 \to 1$ & $0$ 最小的未拜訪子節點是 $1$ \\
3 & $0 \to 1 \to 2$ & $1$ 的子節點 $0$ 已在路徑上，選 $2$ \\
4 & $0 \to 1 \to 2 \to 3$ & $2$ 最小的子節點是 $3$ \\
5 & $0 \to 1 \to 2 \to 3 \to 4$ & $3$ 的子節點 $1$ 已在路徑上，選 $4$；$4$ 唯一的子節點 $0$ 在路徑上，回溯到 $3$ \\
6 & $0 \to 1 \to 2 \to 3 \to 5$ & 到達目標，記下長度 $4$；回溯到 $3$，$3$ 沒有別的子節點，回溯到 $2$ \\
7 & $0 \to 1 \to 2 \to 4$ & $4$ 的子節點 $0$ 在路徑上，回溯；$2$、$1$ 都沒有別的子節點，回溯到 $0$ \\
8 & $0 \to 2$ & $0$ 的下一個子節點 $2$（$2$ 已不在目前的路徑上） \\
9 & $0 \to 2 \to 3$ & \\
10 & $0 \to 2 \to 3 \to 1$ & $1$ 的子節點 $0$、$2$ 都在路徑上，回溯到 $3$ \\
11 & $0 \to 2 \to 3 \to 4$ & $4$ 的子節點 $0$ 在路徑上，回溯到 $3$ \\
12 & $0 \to 2 \to 3 \to 5$ & 到達目標，記下長度 $3$；回溯到 $2$ \\
13 & $0 \to 2 \to 4$ & $4$ 的子節點 $0$ 在路徑上，回溯；$2$、$0$ 都沒有別的子節點，結束 \\
\bottomrule
\end{tabular}
\end{center}

找到的 $0$ 到 $5$ 的路徑有兩條：$0 \to 1 \to 2 \to 3 \to 5$（長度 $4$）和 $0 \to 2 \to 3 \to 5$（長度 $3$）。 依第 5 步選
較短的，答案是 $0 \to 2 \to 3 \to 5$。

這個做法要列出所有簡單路徑，最壞是指數時間。 邊的權重都是 $1$ 時，從 $0$ 做一次 BFS 就能在 $O(|V| + |E|)$ 內得到
同樣的答案（資結書第~21.3~節）：$0$ 的距離 $0$；$1, 2$ 距離 $1$；$3, 4$ 距離 $2$；$5$ 距離 $3$。
\end{pastsoln}
